% !TeX root = ../../main.tex
\subsection{must と have to}
    \fitblankbf[]{must} $\fallingdotseq$ \fitblankbfshort{have}\fitblankbfshort{to} 「〜 \fitblankbf[]{しなければならない} 」
    \vspace{0.5em}

    must を否定すると must not ( \fitblankbf[]{mustn't} )、have to を否定すると\fitblankbf[]{don't have to}となる。
    \vspace{0.5em}

    ※ have は\fitblankbf[]{一般動詞}であるので、否定には基本的に\fitblankbf[]{don't}を使うが、主語の人称、単複、時制によっては
    \fitblankbf[]{doesn't}や\fitblankbf[]{didn't}を使わなければならない。

    % \begin{enumerate}
    %     \item You must finish your homework today.
    %     \dialogueblank{あなたは今日宿題を終えなければならない。}
    %     \vspace{0.5em}
    %     \item I have to get up early tomorrow.
    %     \dialogueblank{私は明日早く起きなければならない。}
    %     \vspace{0.5em}
    %     \item She doesn't have to come to school on Sunday.
    %     \dialogueblank{彼女は日曜日に学校に来る必要はない。}
    % \end{enumerate}

\subsection{助動詞 should}
    should 「〜 \fitblankbf[]{すべきである} 」
    \vspace{0.5em}

    ※不定詞(\fitblankbf[]{to + 動詞の原形})の形容詞的用法で書き換えることもできる。

    % \begin{enumerate}
    %     \item You should eat more vegetables.
    %     \dialogueblank{あなたはもっと野菜を食べるべきだ。}
    %     \vspace{0.5em}
    %     \item We should study harder for the test.
    %     \dialogueblank{私たちはテストのためにもっと一生懸命勉強すべきだ。}
    % \end{enumerate}

\subsection{must 特有の意味}
    must には「しなければならない」以外にも、次のような意味がある。

    \begin{theorembox}{\textbf{must のその他の意味}}
        \begin{enumerate}
            \item \fitblankbf[]{must} は\fitblankbf[]{断定}を表し、「～に \fitblankbf[]{ちがいない} 」という意味になる
            \item \fitblankbf[]{must not} は\fitblankbf[]{禁止}を表し、「～ \fitblankbf[]{してはいけない} 」という意味になる。
            \fitblankbf[]{命令文}で書き換えることができる
            \item don't have to + 動詞の原形は「～ \fitblankbf[]{する必要はない} 」という意味である
        \end{enumerate}
    \end{theorembox}

    % \begin{enumerate}
    %     \item He must be tired.
    %     \dialogueblank{彼は疲れているにちがいない。}
    %     \vspace{0.5em}
    %     \item You must not run in the hallway.
    %     \dialogueblank{廊下を走ってはいけない。}
    %     \vspace{0.5em}
    %     \item You don't have to bring your textbook today.
    %     \dialogueblank{今日は教科書を持ってくる必要はない。}
    % \end{enumerate}

\newpage
